%%%%%%%%%%%%%%%%%%%%%%%%%%%%%%%%%%%%%%%%%%%%%%%%%%%%%%%%%%%%%%%%%%%%
\documentclass[journal]{IEEEtran}
\usepackage[table]{xcolor}
\usepackage{circuitikz}
\usetikzlibrary{calc}
\usepackage{amsmath}
\usepackage{pgfplots}
\pgfplotsset{compat=1.18}
\usepackage{tikz}
\usetikzlibrary{arrows.meta}
\usepackage{graphicx}
\usepackage{booktabs}
\usepackage[utf8]{inputenc}
\usepackage[T1]{fontenc}
\usepackage[backend=biber,style=ieee]{biblatex}
\addbibresource{literatura.bib}
\usepackage[spanish]{babel}
\usepackage{caption}
\usepackage{float}
\usepackage{array}

\hyphenation{op-tical net-works semi-conduc-tor}
\addto\captionsspanish{\renewcommand{\tablename}{Tabla}}
\captionsetup[figure]{justification=centering}
\captionsetup[table]{position=bottom, font=small}
\graphicspath{{./}{./fig/}}

\begin{document}

\title{Entregable 1 -- Proyecto: Semáforo vehicular y peatonal}

\author{Kevin~Méndez~Carvajal,~\IEEEmembership{Estudiante de Ingeniería Electrónica,~ITCR}
        ~Wilson~Osejo~Ruiz,~\IEEEmembership{Estudiante de Ingeniería Electrónica,~ITCR}
        ~Kevin~Sobalvarro~Mora,~\IEEEmembership{Estudiante de Ingeniería Electrónica,~ITCR}
        y~Brandon~Nájera~Zúñiga,~\IEEEmembership{Estudiante de Ingeniería Electrónica,~ITCR}}

\markboth{EL3215 Laboratorio de Electrónica Analógica, IS2026}%
{EL3215 Laboratorio de Electrónica Analógica}

\maketitle

% ---------------------------------------------------------------------------
\begin{abstract}
Espacio para el resumen.
\end{abstract}

\begin{IEEEkeywords}
Semáforo vehicular, semáforo peatonal, temporizador 555, contador BCD, CD4510BE,
HEF4511BP, MC14538BCP, TC4013BP, TC4093BP, display de siete segmentos, amplificador
de audio, regulador de tensión, lógica digital discreta.
\end{IEEEkeywords}

% ===========================================================================
\section{Introducción}

\IEEEPARstart{L}{os} sistemas de control de tráfico son fundamentales para la regulación
segura y eficiente de la movilidad vehicular y peatonal en entornos urbanos. Estos sistemas
integran elementos de temporización, señalización visual y auditiva, así como mecanismos
de interacción con los usuarios~\cite{Floyd2008}.

En el contexto de la ingeniería electrónica, el diseño de un sistema de semáforo permite
aplicar conceptos relacionados con la generación de señales, temporización, control lógico,
visualización de información y procesamiento de señales de audio. Además, representa un
ejemplo claro de integración de subsistemas analógicos y digitales en una aplicación real.

En este proyecto se diseña e implementa un sistema completo de semáforo vehicular y
peatonal utilizando únicamente circuitos digitales discretos —temporizadores, contadores,
compuertas y registros— sin el uso de microcontroladores ni FPGAs, tal como lo exige la
especificación técnica del curso. El sistema es alimentado desde $120\,\text{VAC}$
mediante una fuente diseñada con transformador, rectificador y regulador de tensión
propios.

El sistema se divide en cuatro subsistemas independientes: semáforo vehicular, semáforo
peatonal con display de cuenta regresiva de dos dígitos, sistema de audio para la
señalización auditiva del paso peatonal, y sistema de potencia. Cada subsistema fue
asignado a un integrante del equipo y se describe de forma independiente en las secciones
de marco teórico y metodología.

% ===========================================================================
\section{Marco Teórico}

% ---------------------------------------------------------------------------
\subsection{Semáforo vehicular}
Sección a cargo de: Wilson
\subsection{CD4013 – Flip-flop D dual}

El CD4013 es un circuito integrado CMOS que contiene dos flip-flops tipo D independientes, disparables por flanco de subida del reloj. Cada flip-flop posee entradas asíncronas \texttt{SET} (preset) y \texttt{RESET} (clear) activas en nivel alto, además de salidas directa (\texttt{Q}) y complementada (\texttt{$\overline{Q}$}). Su tabla de verdad establece que en el flanco de subida de \texttt{CLK}, el valor presente en \texttt{D} se transfiere a \texttt{Q}. Las entradas asíncronas tienen prioridad sobre el reloj. Opera con tensiones de alimentación típicas entre 3 V y 15 V, siendo 5 V un valor común en sistemas digitales de baja potencia.

\subsection{LM555CN – Temporizador analógico}

El LM555 es un temporizador integrado de 8 pines que puede configurarse en tres modos principales: astable (oscilador libre), monoestable (generador de pulso único) y biestable (disparador Schmitt). Internamente contiene dos comparadores, un flip-flop RS, un transistor de descarga (\texttt{DIS}) y una etapa de salida capaz de suministrar o disipar hasta 200 mA.

\begin{itemize}
    \item \textbf{Modo monoestable}: Un pulso negativo en \texttt{TRI} (\textit{trigger}) lleva la salida \texttt{OUT} a nivel alto durante un tiempo \(T = 1.1 \cdot R \cdot C\), donde \(R\) es la resistencia entre \texttt{VCC} y \texttt{DIS}, y \(C\) el capacitor entre \texttt{THR} y tierra. Al finalizar, la salida vuelve a baja.
    \item \textbf{Modo astable}: Genera una onda cuadrada continua. La frecuencia y el ciclo de trabajo dependen de dos resistencias (\(R_A\), \(R_B\)) y un capacitor \(C\) conectados a \texttt{DIS}, \texttt{THR} y \texttt{TRI}.
    \item \textbf{Pin \texttt{CON}} (control): permite modificar la tensión de referencia interna de los comparadores; típicamente se desacopla a tierra con un capacitor de 10 nF.
\end{itemize}

\subsection{Transistores bipolares NPN y PNP}

Los transistores NPN (ej. 2N3904, BC547) se usan como interruptores de baja lado: cuando la base recibe corriente suficiente (a través de una resistencia limitadora), el transistor satura y conecta el colector a tierra. Los transistores PNP (ej. 2N3906, BC557) funcionan como interruptores de alto lado: con la base a un nivel bajo respecto del emisor, saturan y conectan el colector a la tensión de alimentación.

% ---------------------------------------------------------------------------
\subsection{Semáforo peatonal y display de cuenta regresiva}
Sección a cargdo de: Brandon
\subsubsection{Oscilador de base de tiempo — NE555 astable}

El circuito integrado NE555, operado en modo astable, constituye la base de tiempo del
subsistema peatonal. En esta configuración el chip oscila de forma autónoma, produciendo
en su salida (pin~3) un tren de pulsos cuadrados cuya frecuencia depende de dos
resistencias externas $R_A$ y $R_B$, y de un capacitor de temporización
$C$~\cite{Floyd2008}:

\begin{equation}
    f = \frac{1{,}44}{(R_A + 2\,R_B)\,C}
    \label{eq:555_freq}
\end{equation}

Con $R_A = 68\,\text{k}\Omega$, $R_B = 47\,\text{k}\Omega$ y $C = 10\,\mu\text{F}$
se obtiene un período de aproximadamente $1{,}1\,\text{s}$ por pulso, suficiente para
implementar la cuenta regresiva de 20 segundos. Los tiempos en alto y en bajo son:

\begin{equation}
    t_H = 0{,}693\,(R_A + R_B)\,C, \qquad t_L = 0{,}693\,R_B\,C
    \label{eq:555_tiempos}
\end{equation}

\subsubsection{Monoestable — MC14538BCP}

El MC14538BCP contiene dos monoestables retrigerables de precisión. Un monoestable
permanece en su estado estable hasta recibir un disparo; entonces su salida cambia
durante un tiempo $T_W$ fijado por componentes externos, y luego regresa al
reposo:

\begin{equation}
    T_W = R_{ext}\,C_{ext}
    \label{eq:mono}
\end{equation}

Cada unidad puede dispararse por flanco positivo ($A_2$) o negativo ($A_1$). En el
sistema peatonal se usan dos instancias.

\paragraph{Primer monoestable — antirrebote del botón.}
Los pulsadores mecánicos producen rebotes de hasta $10\,\text{ms}$. El primer
monoestable se dispara por flanco negativo en $A_1$ cuando el peatón presiona el botón
(la entrada cae de $+5\,\text{V}$ a $0\,\text{V}$ gracias a una resistencia de
pull-up de $10\,\text{k}\Omega$), generando un único pulso limpio de:

\begin{equation}
    T_{W1} = 10\,\text{k}\Omega \times 10\,\mu\text{F} = 100\,\text{ms}
    \label{eq:tw1}
\end{equation}

Este pulso supera el tiempo de rebote típico y actúa como orden puntual hacia el
flip-flop~\cite{Floyd2008}.

\paragraph{Segundo monoestable — señal LOAD.}
Los contadores CD4510BE requieren un pulso en su pin \textit{Preset Enable} ($PE$,
activo alto) para cargar el valor inicial 20. El segundo monoestable detecta el flanco
positivo de $Q$ del flip-flop y genera:

\begin{equation}
    T_{W2} = 1\,\text{k}\Omega \times 1\,\text{nF} = 1\,\mu\text{s}
    \label{eq:tw2}
\end{equation}

\subsubsection{Flip-flop tipo D — TC4013BP}

El TC4013BP contiene dos flip-flops tipo~D con entradas asíncronas de set
($\overline{SD}$) y reset ($\overline{CD}$). La Tabla~\ref{tab:verdad_4013} resume
el comportamiento relevante.

\begin{table}[htbp]
    \centering
    \caption{Entradas asíncronas del TC4013BP}
    \label{tab:verdad_4013}
    \begin{tabular}{ccc}
        \toprule
        $\overline{SD}$ & $\overline{CD}$ & $Q$ \\
        \midrule
        1 & 0 & 1 (set) \\
        0 & 1 & 0 (reset) \\
        0 & 0 & indefinido \\
        \bottomrule
    \end{tabular}
\end{table}

El pulso del primer monoestable activa $\overline{SD}$ ($Q \to 1$, paso habilitado).
Al finalizar la cuenta, la lógica de detección activa $\overline{CD}$ ($Q \to 0$,
reposo). Las salidas $Q$ y $\overline{Q}$ controlan los LED verde y rojo,
respectivamente, y el pin de blanking de los decodificadores.

\subsubsection{Contadores BCD descendentes — CD4510BE}

El CD4510BE es un contador BCD síncrono de 4 bits con dirección seleccionable
($U/\overline{D} = 0$ para descender), habilitación por $CI$ y carga paralela asíncrona
por $PE$. Para representar 20 se usan dos contadores en cascada: unidades con
$P_1$--$P_4 = 0000$ (BCD~0) y decenas con $P_2 = 1$, demás a tierra (BCD~2). El $CO$
del contador de unidades habilita el conteo de las decenas:

\begin{equation}
    N_{total} = N_{dec} \times 10 + N_{uni}
    \label{eq:bcd}
\end{equation}

Cuando ambos llegan a cero, el $CO$ del contador de decenas activa la lógica de reset.

\subsubsection{Decodificador BCD a siete segmentos — HEF4511BP}

El HEF4511BP convierte el BCD de 4 bits en las señales $a$--$g$ para un display de
cátodo común. El pin $\overline{BI}$ (blanking, activo bajo) se conecta a $Q$ del
flip-flop para apagar el display en estado de reposo. La corriente por segmento se
limita con $R_{seg} = 220\,\Omega$:

\begin{equation}
    I_{seg} = \frac{V_{OH} - V_F}{R_{seg}} \approx \frac{4{,}5 - 2{,}0}{220}
            \approx 11{,}4\,\text{mA}
    \label{eq:iseg}
\end{equation}

\subsubsection{Oscilador de base de tiempo — NE555 astable}
El circuito integrado NE555, operado en modo astable, constituye la base de tiempo del
subsistema peatonal. ... (se mantiene igual)

% [El resto de subsecciones se mantienen igual hasta los componentes]

\subsubsection{Componentes pasivos del subsistema peatonal}

Las Tablas~\ref{tab:res} y \ref{tab:cap} listan las resistencias y capacitores
utilizados en este subsistema.

\begin{table}[htbp]
    \centering
    \caption{Resistencias del subsistema peatonal}
    \label{tab:res}
    \begin{tabular}{lcc}
        \toprule
        \textbf{Referencia} & \textbf{Valor nominal} & \textbf{Función} \\
        \midrule
        $R_1$            & $68\,\text{k}\Omega$   & $R_A$ del 555 astable \\
        $R_2$            & $47\,\text{k}\Omega$   & $R_B$ del 555 astable \\
        $R_3$, $R_4$     & $10\,\text{k}\Omega$   & Pull-up botones A y B \\
        $R_5$            & $10\,\text{k}\Omega$   & $R_{ext}$ monoestable 1 (antirrebote) \\
        $R_6$            & $1\,\text{k}\Omega$    & $R_{ext}$ monoestable 2 (LOAD) \\
        $R_7$--$R_{20}$  & $330\,\Omega$          & Limitación de corriente \\
        $R_{21}$--$R_{24}$ & $330\,\Omega$        & Limitación LED verde y rojo \\
        \bottomrule
    \end{tabular}
\end{table}

\begin{table}[htbp]
    \centering
    \caption{Capacitores del subsistema peatonal}
    \label{tab:cap}
    \begin{tabular}{lcc}
        \toprule
        \textbf{Referencia} & \textbf{Valor nominal} & \textbf{Función} \\
        \midrule
        $C_1$        & $10\,\mu\text{F}$     & Temporización del oscilador 555 \\
        $C_2$        & $10\,\mu\text{F}$     & $C_{ext}$ monoestable 1 (antirrebote) \\
        $C_3$        & $1\,\text{nF}$        & $C_{ext}$ monoestable 2 (LOAD) \\
        $C_4$        & $10\,\text{nF}$       & Desacoplo / filtrado adicional \\
        $C_5$        & $0{,}01\,\text{nF}$   & Desacoplo de alta frecuencia \\
        $C_6$, $C_7$ & $100\,\text{nF}$      & Desacoplo de alimentación (bypass) \\
        \bottomrule
    \end{tabular}
\end{table}

% ---------------------------------------------------------------------------
\subsection{Sistema de audio}
Sección a cargo de: Kevin Méndez
TC4066BP:
El TC4066BP contiene cuatro interruptores analógicos bidireccionales controlados digitalmente. Cada interruptor posee tres terminales: dos de señal y uno de control. Cuando el pin de control recibe un nivel lógico alto (+5V), el canal entre los terminales de señal conduce con una resistencia típica de 80 Ω, despreciable frente a los componentes externos del circuito. Cuando el control recibe un nivel bajo (0V), el canal se abre presentando impedancia virtualmente infinita.
A diferencia de los interruptores mecánicos, la conmutación es instantánea y no genera rebotes, lo que lo hace ideal para seleccionar componentes pasivos en circuitos de temporización y oscilación Floyd2008.

% ---------------------------------------------------------------------------
\subsection{Sistema de potencia y transformador}
% Sección a cargo de: (nombre del compañero)
Texto.

% ===========================================================================
\section{Metodología}

% ---------------------------------------------------------------------------
\subsection{Semáforo vehicular}
Sección a cargo de: Wilson
\subsection{Empleo del CD4013}

El flip-flop se utiliza como elemento de memoria del estado principal del semáforo: \texttt{1Q} alto indica que el LED verde está encendido y el sistema en fase ``verde''; \texttt{1Q} bajo indica fase ``rojo/amarillo''. 

\begin{itemize}
    \item La entrada \texttt{1D} se fija a VCC, de modo que cada flanco de subida en \texttt{1CLK} (generado por un pulsador) llevará \texttt{1Q} a nivel alto. 
    \item La entrada asíncrona \texttt{1RST} se conecta al colector de un transistor NPN (Q4) controlado por la salida del monoestable del LED rojo (CI2). Cuando el LED rojo termina su temporización, Q4 satura, llevando \texttt{1RST} a nivel bajo, lo que resetea el flip-flop y pone \texttt{1Q} en bajo, finalizando la fase verde.
    \item Las demás entradas (\texttt{1SET}, \texttt{2SET}, \texttt{2D}, \texttt{2CLK}, \texttt{2RST}) se conectan a tierra para inutilizar el segundo flip-flop.
\end{itemize}

\subsection{Empleo de los tres LM555}

\subsubsection{CI4 (555 \#1) – Oscilador astable para parpadeo del LED amarillo}

Este temporizador se configura en modo astable libre con dos resistencias de 4.7 k$\Omega$ (R12 y R13) y un capacitor de 47 $\mu$F (C5). Genera una onda cuadrada en su salida \texttt{OUT} que se aplica al ánodo del LED amarillo a través de una resistencia de 500 $\Omega$. La frecuencia de parpadeo permite que el LED amarillo destelle cuando está habilitado.

\subsubsection{CI3 (555 \#2) – Monoestable para temporización del LED amarillo y disparo del rojo}

Este circuito actúa como generador de pulso único. Su entrada \texttt{TRI} está normalmente en alto (por resistencia R9 de 10 k$\Omega$ a VCC). Cuando \texttt{1Q} del 4013 se pone a alto, el transistor Q1 satura y lleva \texttt{TRI} a nivel bajo, disparando el monoestable. 

La salida \texttt{OUT} de CI3 pasa entonces a nivel alto durante un tiempo determinado por R8 (470 k$\Omega$) y C3 (10 $\mu$F). Esta salida cumple dos funciones:
\begin{enumerate}
    \item A través de Q3 (NPN), conecta el cátodo del LED amarillo a tierra, permitiendo que el LED amarillo se ilumine con la señal de parpadeo proveniente de CI4.
    \item Mediante una red diferenciadora (C7 de 10 $\mu$F + R11 de 10 k$\Omega$), genera un pulso corto negativo en el nodo de unión que se conecta al pin \texttt{TRI} de CI2 (555 \#3), iniciando así la temporización del LED rojo.
\end{enumerate}

Al finalizar el tiempo del monoestable, \texttt{OUT} de CI3 vuelve a nivel bajo, apagando el LED amarillo y deshabilitando el paso de la señal de CI4.

\subsubsection{CI2 (555 \#3) – Monoestable para el LED rojo y reset del sistema}

Este temporizador también funciona como monoestable. Su entrada \texttt{TRI} recibe el pulso generado por la red diferenciadora de CI3. Al dispararse, su salida \texttt{OUT} se pone a nivel alto durante un tiempo nominal dado por R5 (1.8 M$\Omega$) y C1 (10 $\mu$F). 

La salida alta:
\begin{itemize}
    \item Enciende el LED rojo a través de R6 (330 $\Omega$).
    \item Satura el transistor Q4 mediante R7 (10 k$\Omega$). El colector de Q4 se conecta a \texttt{1RST} del 4013 a través de R4 (10 k$\Omega$), forzando el reset del flip-flop y por tanto apagando el LED verde.
\end{itemize}
Al terminar la temporización, \texttt{OUT} de CI2 vuelve a nivel bajo, el LED rojo se apaga y Q4 deja de saturar, liberando \texttt{1RST} y dejando el sistema listo para un nuevo ciclo al presionar nuevamente el botón.

\subsection{Control de los LEDs}

\begin{itemize}
    \item \textbf{LED verde}: Controlado por el transistor PNP Q2. Cuando \texttt{1Q} del 4013 es alto, la base de Q2 se lleva a un nivel bajo (respecto del emisor a VCC) a través de R2 (10 $\Omega$), saturando el PNP y permitiendo la corriente desde VCC al LED verde limitada por R3 (660 $\Omega$).
    \item \textbf{LED amarillo}: El ánodo recibe la señal de parpadeo desde CI4. El cátodo es conectado a tierra por Q3 solo cuando la salida de CI3 es alta. De este modo, el LED amarillo parpadea únicamente durante el tiempo que dura el pulso del monoestable CI3.
    \item \textbf{LED rojo}: Se enciende directamente por la salida de CI2 cuando su monoestable está activo.
\end{itemize}

\subsection{Diagrama de flujo de la secuencia}

1. Reposo inicial: 4013 reseteado, LED rojo apagado (si CI2 no está disparado), LED amarillo apagado, LED verde apagado.
2. Usuario presiona botón → flanco en \texttt{1CLK} → \texttt{1Q} = alto → LED verde encendido.
3. \texttt{1Q} alto → Q1 satura → dispara CI3 (555\#2).
4. CI3 genera pulso de tiempo fijo:
   - Durante ese pulso: LED amarillo parpadea (por CI4 + Q3).
   - Al final del pulso, la red RC (C7+R11) dispara CI2.
5. CI2 genera pulso de tiempo fijo:
   - Durante ese pulso: LED rojo encendido.
   - Q4 saturado → mantiene \texttt{1RST} del 4013 en bajo → \texttt{1Q} = bajo → LED verde apagado.
6. Termina pulso de CI2 → LED rojo apagado, Q4 corta → \texttt{1RST} liberado.
7. Sistema vuelve al estado 1, esperando nuevo pulso de botón.

\subsection{Componentes discretos utilizados}

Los valores de resistencias y capacitores se detallan en la Tabla~\ref{tab:pasivos} (incluida en el anexo correspondiente). Todos los capacitores electrolíticos respetan su polaridad, y los capacitores cerámicos de 10 nF se emplean para desacoplo del pin \texttt{CON} de cada 555.
% ---------------------------------------------------------------------------
\subsection{Semáforo peatonal y display de cuenta regresiva}
A cargo de: Brandon

El subsistema peatonal opera de la siguiente forma: en estado de reposo, ambos semáforos
muestran luz roja y los displays permanecen apagados ($Q = 0$, $\overline{BI}$ activo).
Al presionar cualquiera de los dos botones, el primer monoestable genera un pulso limpio
que activa el flip-flop TC4013BP ($Q \to 1$). En ese instante el segundo monoestable
detecta el flanco positivo y genera el pulso LOAD, cargando el valor 20 en los contadores.
El oscilador de $\approx 1\,\text{Hz}$ decrementa los contadores segundo a segundo. Al
llegar a cero, el TC4093BP invierte el $CO$ de las decenas y resetea el flip-flop,
apagando el verde y devolviendo el sistema al estado de reposo.

La señal de solicitud peatonal ($Q$) se comunica al subsistema vehicular para forzar su
transición a rojo.

\subsubsection{Observaciones sobre la simulación en NI Multisim}

Durante la etapa de simulación se identificaron limitaciones que impidieron verificar el
funcionamiento integral del subsistema. Los principales problemas encontrados fueron:

\begin{itemize}
    \item Los pines ocultos de alimentación (VDD y VSS) de los circuitos integrados CMOS
          no se conectan automáticamente; si los símbolos globales no se colocan de forma
          explícita, los chips quedan sin tensión y sus salidas son indeterminadas.
    \item El pin $CI$ del contador de decenas no quedó correctamente enlazado al $CO$ del
          contador de unidades en algunas versiones del esquema, impidiendo el decremento
          de las decenas.
\end{itemize}

% ---------------------------------------------------------------------------
\subsection{Sistema de audio}
A cargo de: Kevin Méndez

El subsistema de audio opera de la siguiente forma: en estado de reposo, 
con el semáforo vehicular en verde, el flip-flop del subsistema peatonal 
permanece inactivo ($Q = 0$), manteniendo el NE555 inhibido mediante su 
pin de reset (pin 4) y silenciando el parlante. Al activarse el cruce 
peatonal ($Q \rightarrow 1$), el NE555 comienza a oscilar a 1 Hz, 
produciendo el pitido característico a través del TC4066BP, que conecta 
la resistencia $R_{B1} = 68\,k\Omega$ al pin de descarga del oscilador. 
Simultáneamente, el MC14538BCP inicia una temporización de 15 segundos 
($T_W = R_{ext}C_{ext}$). Al concluir este intervalo, su flanco de salida 
setea el TC4013BP ($Q \rightarrow 1$), el cual conmuta el TC4066BP hacia 
la resistencia $R_{B2} = 33\,k\Omega$, elevando la frecuencia de 
oscilación a 2 Hz durante los últimos 5 segundos del cruce. Al finalizar 
el verde peatonal, la señal $Q$ del flip-flop principal regresa a 0, 
reseteando el TC4013BP y el NE555, devolviendo el sistema al estado de 
reposo.

% ---------------------------------------------------------------------------
\subsection{Sistema de potencia y transformador}
% Sección a cargo de: (nombre del compañero)
Texto.

% ===========================================================================
\section{Análisis de resultados}
Texto.

% ===========================================================================
\section{Conclusiones y Recomendaciones}

\subsection{Conclusiones}
Texto.

% ===========================================================================
\nocite{*}
\printbibliography

% ===========================================================================
\begin{IEEEbiographynophoto}{Kevin Méndez Carvajal}
Kevin Méndez Carvajal es estudiante de Ingeniería Electrónica en el Tecnológico de
Costa Rica (ITCR). Su área de interés incluye el diseño de circuitos analógicos,
sistemas embebidos y automatización industrial.
\end{IEEEbiographynophoto}

\begin{IEEEbiographynophoto}{Wilson Osejo Ruiz}
Wilson Osejo Ruiz es estudiante de Ingeniería Electrónica en el Tecnológico de
Costa Rica (ITCR). Sus intereses se centran en electrónica de potencia,
instrumentación y control de procesos.
\end{IEEEbiographynophoto}

\begin{IEEEbiographynophoto}{Kevin~Sobalvarro~Mora}
Kevin Sobalvarro Mora es estudiante de Ingeniería Electrónica en el Tecnológico de
Costa Rica (ITCR). Cursando el curso EL3215 Laboratorio de Electrónica Analógica,
I Semestre 2026.

Correo: kevinmorasoval03@estudiantec.cr
\end{IEEEbiographynophoto}

\begin{IEEEbiographynophoto}{Brandon~Nájera~Zúñiga}
Brandon Nájera Zúñiga es estudiante de Ingeniería Electrónica en el Tecnológico de
Costa Rica (ITCR). Cursando el curso EL3215 Laboratorio de Electrónica Analógica,
I Semestre 2026.

Correo: brandonnz@estudiantec.cr
\end{IEEEbiographynophoto}

\vfill

\end{document}